\title{DeepSeekMath: Pushing the Limits of Mathematical Reasoning in Open Language Models}

\begin{abstract}
Language models often struggle with mathematical reasoning, primarily due to the intricate and highly structured characteristics inherent in mathematics. This work presents DeepSeekMath~7B, an extension of DeepSeek-Coder-Base-v1.5 7B, which undergoes continued pre-training on 120B math-centric tokens extracted from Common Crawl, as well as additional natural language and programming data. DeepSeekMath~7B attains a remarkable 51.7\% score on the competition-level MATH benchmark—without the aid of external toolkits or ensembling methods—placing its performance on par with leading systems such as Gemini-Ultra and GPT-4. Notably, leveraging self-consistency over 64 generated samples, DeepSeekMath~7B achieves 60.9\% accuracy on MATH. Two principal elements underpin the model’s proficiency in mathematical reasoning: First, the effective utilization of openly available web data, curated through a carefully designed selection process; second, the implementation of Group Relative Policy Optimization (GRPO), an adaptation of Proximal Policy Optimization (PPO), which simultaneously boosts mathematical reasoning ability and optimizes PPO’s memory efficiency.
\end{abstract}

\section{Introduction}

The advent of large language models (LLMs) has transformed mathematical reasoning within artificial intelligence, propelling notable progress in both quantitative \citep{MATH} and geometry \citep{trinh2024solving} reasoning benchmarks. Additionally, these systems have emerged as valuable tools for supporting humans in tackling intricate mathematical challenges \citep{tao}. Despite these advances, top-performing models like GPT-4 \citep{gpt4} and Gemini-Ultra \citep{gemini} remain proprietary, while the open-source models that are publicly available continue to fall substantially short in terms of mathematical performance.

This work presents DeepSeekMath, a specialized language model for mathematics that achieves superior results compared to existing open-source alternatives and attains performance near that of GPT-4 across academic evaluation tasks. Central to our approach is the construction of the DeepSeekMath Corpus, a large-scale, high-quality pre-training dataset consisting of 120B mathematical tokens. This dataset is derived from Common Crawl (CC) with the assistance of a classifier built on fastText \citep{joulin2016fasttext}. In the classifier’s initial training phase, we leverage positive samples from OpenWebMath \citep{openwebmath} alongside a heterogeneous collection of unrelated webpages as negative samples. The trained classifier is then utilized to identify additional mathematical content within the CC; these positive samples are subjected to further screening and refinement by human annotators. Using this expanded and improved dataset, the classifier undergoes a subsequent training iteration to further enhance its discriminatory power. Our evaluation demonstrates the corpus’s strong quality: the DeepSeekMath-Base 7B model attains 64.2\% accuracy on GSM8K \citep{gsm8k} and 36.2\% on the MATH competition benchmark \citep{MATH}, both surpassing the performance of Minerva 540B \citep{minerva}. Additionally, thanks to the corpus’s multilingual nature, we observe gains on Chinese mathematical assessment tasks \citep{wei2023cmath,agieval}. We view our methodologies for mathematical data preparation as an initial step for the field, highlighting opportunities for substantial progress ahead.

DeepSeekMath-Base is constructed upon DeepSeek-Coder-Base-v1.5 7B \citep{deepseek-coder}, as preliminary results indicate that models trained on code provide a superior foundation compared to those trained on more general language tasks. Notably, training on mathematical data not only strengthens mathematical proficiency but also leads to measurable gains on general reasoning benchmarks such as MMLU \citep{mmlu} and BBH \citep{bbh}.

Subsequent to pre-training, DeepSeekMath-Base undergoes mathematical instruction fine-tuning using datasets focused on chain-of-thought \citep{cot}, program-of-thought \citep{pot,pal}, and tool-augmented reasoning \citep{tora}. The resulting DeepSeekMath-Instruct 7B surpasses all other 7B models, while delivering competitive performance on par with instruction-tuned open-source models of the 70B scale.

Additionally, we present Group Relative Policy Optimization (GRPO), a new reinforcement learning (RL) algorithm based on Proximal Policy Optimization (PPO) \citep{schulman2017proximal}. Unlike PPO, GRPO eliminates the need for a separate critic model by inferring the baseline from group-level scores, greatly lowering the resource requirements for training. Using only a selection of English instruction-tuning data, GRPO achieves marked improvements over the already strong DeepSeekMath-Instruct, boosting accuracy on in-domain tasks (GSM8K: 82.9\% $\rightarrow$ 88.2\%, MATH: 46.8\% $\rightarrow$ 51.7\%) and enhancing out-of-domain generalization (e.g., CMATH: 84.6\% $\rightarrow$ 88.8\%) during RL training. We also establish a unified conceptual framework that relates a variety of optimization strategies—including Rejection Sampling Fine-Tuning (RFT) \citep{yuan2023scaling}, Direct Preference Optimization (DPO) \citep{dpo}, PPO, and GRPO—showing that these approaches fall under the umbrella of direct or simplified RL methodologies. Through comprehensive experimentation—including comparisons of online versus offline updates, result- versus process-focused supervision, and single-step versus iterative RL—we analyze critical aspects of this paradigm. Finally, we discuss the mechanisms by which our RL method enhances instruction-tuned model performance, and highlight future research paths for more efficient RL within the scope of this unified approach.

\subsection{Contributions}
This work presents advancements in large-scale mathematical pre-training as well as an in-depth investigation of reinforcement learning techniques.

\textbf{Math Pre-Training at Scale}

\begin{itemize}[topsep=0pt]
    \item We demonstrate that the Common Crawl public dataset is a rich resource for mathematics-related data. By employing a rigorous data selection process, we have curated the DeepSeekMath Corpus, comprising 120B tokens of mathematically-relevant web text—a dataset nearly 7 times larger than the math web corpus used by Minerva \citep{minerva} and 9 times greater than the newly introduced OpenWebMath \citep{openwebmath}.
    
    \item Our DeepSeekMath-Base 7B base model, pre-trained on this corpus, attains performance on par with Minerva 540B \citep{minerva}, suggesting that sheer parameter count does not exclusively determine proficiency in mathematical reasoning. Substantially smaller models, when pre-trained on high-quality mathematical data, can perform equally well.
    
    \item We also report insights gleaned from our pre-training experiments. Models subjected to code-focused training before mathematical instruction exhibit improved problem-solving capabilities in mathematics, both when leveraging tools and operating independently. This finding provides partial evidence in addressing the persistent query: \textit{does training on code facilitate better reasoning?} Our results affirm that, particularly in mathematical domains, it indeed does.
    
    \item Despite the prevalence of arXiv-based training in math-related research, our empirical evaluations reveal that incorporating arXiv papers does not yield significant performance gains across the suite of mathematical benchmarks employed in this study.
\end{itemize}

\textbf{Exploration and Analysis of Reinforcement Learning}

\begin{itemize}[topsep=0pt]
    \item We present Group Relative Policy Optimization (GRPO), a reinforcement learning framework that is both efficient and high-performing. GRPO dispenses with the traditional critic component and instead derives the baseline from aggregated group scores, leading to notable reductions in computational requirements during training when compared to Proximal Policy Optimization (PPO).
    \item Our results show that GRPO markedly improves the instruction-tuned DeepSeekMath-Instruct model by leveraging only instruction-tuning datasets. Moreover, we notice improved generalization to out-of-domain tasks throughout the reinforcement learning phase.
    \item We offer a comprehensive framework that clarifies the relationships among various methods, including RFT, DPO, PPO, and GRPO. Additionally, we carry out a wide array of experiments—such as online versus offline training, outcome-based versus process-based supervision, and single-turn versus iterative reinforcement learning—to thoroughly examine the core factors of this unified paradigm.
    \item Leveraging this framework, we investigate why reinforcement learning is successful for LLMs and outline several promising avenues for enhancing reinforcement learning approaches in this context.
\end{itemize}

\subsection{Summary of Evaluations and Metrics}

\begin{itemize}[topsep=0pt]
    \item \textbf{English and Chinese Mathematical Reasoning}: We thoroughly evaluate our models across diverse English and Chinese mathematical benchmarks, spanning tasks from elementary to collegiate complexity. The English datasets assessed include GSM8K \citep{gsm8k}, MATH \citep{MATH}, SAT \citep{llemma}, OCW Courses \citep{minerva}, and MMLU-STEM \citep{mmlu}. For Chinese, we utilize MGSM-zh \citep{mgsm}, CMATH \citep{wei2023cmath}, Gaokao-MathCloze \citep{agieval}, and Gaokao-MathQA \citep{agieval}. Evaluation criteria encompass both the capacity for the models to autonomously generate textual problem solutions and their proficiency in employing Python to resolve questions.

    On English tasks, DeepSeekMath-Base exhibits performance comparable to that of the proprietary Minerva 540B model \citep{minerva}, and clearly outperforms all available open-source base models (such as Mistral 7B \citep{mistral} and Llemma-34B \citep{llemma}), with or without mathematical pre-training—frequently by a considerable margin. On Chinese assessments, DeepSeekMath-Base excels as well, a result we attribute to the inclusion of high-quality multilingual pre-training data, diverging from the monolingual English emphasis observed in prior work \citep{minerva,llemma}. With the addition of instruction tuning and reinforcement learning, the DeepSeekMath-Instruct and DeepSeekMath-RL models achieve notable results, exceeding 50\% accuracy on the challenging MATH dataset for the first time among open-source models.
\end{itemize}

\item \textbf{Formal Mathematics}: The performance of DeepSeekMath-Base is assessed on the informal-to-formal theorem proving task introduced in \citep{dsp_proof}, utilizing miniF2F \citep{minif2f} with Isabelle \citep{isabelle} serving as the proof assistant. Results indicate that DeepSeekMath-Base achieves notable success in few-shot autoformalization.

\item \textbf{Natural Language Understanding, Reasoning, and Code}: To thoroughly evaluate the model’s proficiency in general language comprehension, reasoning, and programming, we test DeepSeekMath-Base on several benchmarks. These include the Massive Multitask Language Understanding (MMLU) benchmark \citep{mmlu}, consisting of 57 multiple-choice tasks from diverse subject areas; BIG-Bench Hard (BBH) \citep{bbh}, which comprises 23 intricate tasks that generally require multi-step reasoning; and two widely recognized code evaluation benchmarks, HumanEval \citep{codex} and MBPP \citep{mbpp}. Mathematical pre-training demonstrably boosts both the model’s linguistic understanding and its reasoning abilities.
\end{itemize}

\section{Math Pre-Training}

\subsection{Data Collection and Decontamination} This section details the creation of the DeepSeekMath Corpus sourced from Common Crawl. As visualized in Figure \ref{fig:data_collect}, we employ a multi-stage pipeline that efficiently aggregates a large-scale mathematical dataset from Common Crawl, initializing the process with a high-quality, curated seed set of mathematical texts. The same methodology is readily adaptable to alternative fields, including programming.

Initially, we select OpenWebMath \citep{openwebmath}—a curated assembly of mathematical content from the web—as our foundational seed dataset. A fastText model \citep{joulin2016fasttext} is then trained using this seed to identify additional web pages with similar mathematical characteristics. For model training, we sample 500,000 positive examples from the seed and pair them with 500,000 randomly chosen Common Crawl web pages serving as negative samples. Model training is conducted with an open-source implementation\footnote{\url{https://fasttext.cc}}, using hyperparameters: vector dimension of 256, learning rate set to 0.1, maximum word n-gram length of 3, minimum word frequency of 3, and 3 training epochs. In order to scale down the initial size of Common Crawl, both exact and approximate (near) deduplication are carried out at the URL level, reducing the data set to 40B HTML web pages. Mathematical web documents are retrieved from this deduplicated set using the trained fastText model. To ensure high content quality, the retrieved documents are scored according to the model’s predictions, and only the top-ranking documents are retained. The quantity of preserved data is empirically evaluated by pre-training on the leading 40B, 80B, 120B, and 160B tokens. For the first round, we opt to retain the uppermost 40B tokens.

Following the initial round of data collection, a substantial number of mathematical web pages remain uncollected. This is largely due to the limited diversity within the positive samples used to train the fastText model. To enhance the fastText model’s effectiveness, we broaden the seed dataset by identifying supplementary mathematical web domains. Specifically, we partition the entirety of Common Crawl into mutually exclusive domains, with each domain corresponding to web pages sharing the same base URL. For each domain, we determine the proportion of pages that were retrieved during the first iteration. Domains where more than 10\% of the pages were collected are categorized as being math-oriented (for instance, \url{mathoverflow.net}). 

Next, we manually annotate the URLs within these math-related domains that pertain to mathematical content (e.g., \url{mathoverflow.net/questions}). Any previously uncollected web page associated with these annotated URLs is subsequently incorporated into the seed dataset. Through this methodology, we accumulate additional positive examples, which in turn improves the fastText model’s ability to identify mathematical content in later cycles of data collection. After completing four iterations, the dataset expands to 35.5M mathematical web pages, totaling 120B tokens. By the fourth iteration, approximately 98\% of the pages have already been gathered by the third iteration, prompting us to halt further data collection.

In order to prevent benchmark contamination, we adhere to the filtering protocols outlined in \cite{deepseek-coder}, excluding web pages that include questions or answers originating from established English mathematical benchmarks such as GSM8K \citep{gsm8k}, MATH \citep{MATH}, and Chinese benchmarks such as CMATH \citep{wei2023cmath} and AGIEval \citep{agieval}. The filtration process operates as follows: any text fragment containing a 10-gram sequence that perfectly matches a substring from any evaluation benchmark is excluded from the training corpus. For benchmark excerpts shorter than 10 grams but at least 3 grams in length, we also perform exact matching to remove potential overlaps from our collection of mathematical web pages.

\subsection{Validating the Quality of the DeepSeekMath~Corpus}

We conduct pre-training studies to assess how the DeepSeekMath~Corpus performs relative to other contemporary math-focused corpora:

\begin{itemize}[topsep=0pt]
    \item \textbf{MathPile} \citep{mathpile}: a collection spanning multiple sources (8.9B tokens), including textbooks, Wikipedia, ProofWiki, CommonCrawl, StackExchange, and arXiv, with the majority (over 85\%) sourced from arXiv;
    \item \textbf{OpenWebMath} \citep{openwebmath}: a corpus curated from CommonCrawl data filtered specifically for mathematics, comprising 13.6B tokens;
    \item \textbf{Proof-Pile-2} \citep{llemma}: a dataset integrating OpenWebMath, AlgebraicStack (10.3B tokens of mathematical code), and arXiv papers (28.0B tokens). When evaluating on Proof-Pile-2, we adopt the arXiv:Web:Code distribution ratio of 2:4:1 as described in \cite{llemma}.
\end{itemize}

\subsubsection{Training Setting}
\label{sec:quality-policy}
We perform math-specific training on a broadly pre-trained language model comprising 1.3B parameters, architecturally consistent with the DeepSeek LLMs \citep{deepseek-llm}, referred to here as DeepSeek-LLM 1.3B. Each mathematical dataset is used to independently train a separate model, with each trained on a corpus totaling 150B tokens. Training is carried out utilizing the resource-efficient HAI-LLM framework \citep{haillm}. Adhering to the methodology established for DeepSeek LLMs, we optimize the models with AdamW \citep{adamW}, setting $\beta_1=0.9$, $\beta_2=0.95$, and $\mathrm{weight\_decay}=0.1$. The learning rate schedule consists of a warmup phase over 2,000 steps, after which the learning rate is gradually reduced—falling to 31.6\% of its peak by 80\% of training completion and to 10.0\% of the maximum by the 90\% mark. The learning rate is capped at 5.3e-4, while batches comprise 4M tokens and sequences are processed up to 4K tokens in length.

\subsubsection{Evaluation Results}

\textbf{The DeepSeekMath~Corpus distinguishes itself through high quality, extensive multilingual content, and superior scale.}

\begin{itemize}[topsep=0pt]
    \item \textbf{High-quality}:
    We assess downstream capabilities using 8 mathematical benchmark datasets, utilizing few-shot chain-of-thought prompting as in \cite{cot}. Table \ref{tab:corpora_comparison} presents evidence that the model trained on the DeepSeekMath~Corpus significantly outperforms others. Additionally, Figure \ref{fig:corpora_comparisons} demonstrates that after processing 50B tokens (equivalent to one complete epoch of Proof-Pile-2), the DeepSeekMath-trained model surpasses its Proof-Pile-2 counterpart, attesting to the superior average quality of DeepSeekMath.
    \item \textbf{Multilingual}:
    Data within the DeepSeekMath~Corpus is represented in multiple languages, with English and Chinese forming the largest linguistic subsets. Results from Table \ref{tab:corpora_comparison} reveal that pre-training on the DeepSeekMath~Corpus bolsters mathematical reasoning proficiency in both English and Chinese. Conversely, previous mathematical datasets—primarily English-focused—either offer limited benefit or negatively affect Chinese mathematical reasoning outcomes.
    \item \textbf{Large-scale}:
    In terms of volume, the DeepSeekMath~Corpus substantially exceeds the size of prior mathematical corpora. Figure \ref{fig:corpora_comparisons} illustrates that models (DeepSeek-LLM 1.3B) trained on DeepSeekMath~Corpus exhibit accelerated and sustained learning gains. By contrast, models trained on the smaller baseline corpora rapidly exhaust their training data through repeated passes, resulting in stagnated performance improvements.
\end{itemize}

\subsection{Training and Evaluating DeepSeekMath-Base 7B}

This section presents DeepSeekMath-Base 7B, a foundational model that exhibits advanced mathematical reasoning skills. The model leverages initialization from DeepSeek-Coder-Base-v1.5 7B \citep{deepseek-coder} and undergoes training over 500B tokens. The dataset comprises 56\% samples from the DeepSeekMath Corpus, 4\% from AlgebraicStack, 10\% originating from arXiv, 20\% sourced from Github code, and the last 10\% made up of natural language content from Common Crawl in both English and Chinese. The majority of the training protocol aligns with Section \ref{sec:quality-policy}, with two main exceptions: the peak learning rate is set to 4.2e-4, and the token batch size is fixed at 10M.

To thoroughly evaluate the mathematical proficiency of DeepSeekMath-Base 7B, we analyze its capabilities in three aspects: autonomous generation of mathematical solutions without external aids, tool-assisted problem-solving, and formal theorem verification. Additionally, the general aptitude of the base model is explored, highlighting its competencies in natural language understanding, logical reasoning, and coding.

\paragraph{Mathematical Problem Solving with Step-by-Step Reasoning}

We assess the problem-solving effectiveness of DeepSeekMath-Base using few-shot chain-of-thought prompting \citep{cot} on eight separate benchmarks in both English and Chinese. These benchmarks span various types of quantitative reasoning—including GSM8K \citep{gsm8k}, MATH \citep{MATH}, and CMATH \citep{wei2023cmath}—as well as multiple-choice formats such as MMLU-STEM \citep{mmlu} and Gaokao-MathQA \citep{agieval}, thereby reflecting a broad spectrum of mathematical domains from elementary through collegiate levels.

As illustrated in Table \ref{tab:base_math_cot}, DeepSeekMath-Base 7B attains the highest scores across all eight benchmarks among open-source base models, outperforming both the commonly used general model Mistral 7B \citep{mistral} and the newly introduced Llemma 34B \citep{llemma}, which is trained on the Proof-Pile-2 dataset \citep{llemma} for mathematical tasks. In particular, on the challenging MATH competition dataset, DeepSeekMath-Base exceeds the performance of existing open-source base models by a margin greater than 10\% in absolute terms, and it also surpasses the proprietary Minerva 540B \citep{minerva}—a model built upon PaLM \citep{palm} and further refined on mathematical corpora—even though Minerva has 77 times more parameters.

\paragraph{Mathematical Problem Solving with Tool Use} We assess program-based mathematical reasoning using few-shot program-of-thought prompts \citep{pot,pal} on GSM8K and MATH. Each model is given the task of solving problems by composing Python scripts that may utilize libraries such as \textit{math} and \textit{sympy} to handle complex calculations. The final output of the executed script is taken as the response. According to the results in Table \ref{tab:base_math_pot_proof}, DeepSeekMath-Base 7B achieves superior accuracy compared to the previous top-performing Llemma 34B.

\paragraph{Formal Mathematics}

Automating the construction of formal proofs has become increasingly important for guaranteeing the precision and reliability of mathematical arguments, as well as for improving workflow efficiency. We test DeepSeekMath-Base 7B on the informal-to-formal proof generation task described in \citep{dsp_proof}, where the model receives an informal statement, its formal version, and an informal justification, and must produce a corresponding formal proof. Our experiments use the miniF2F \citep{minif2f} benchmark, which targets formal proofs at the Olympiad level; for each challenge, we employ few-shot prompting to generate a formal proof in Isabelle. Following the protocol of \cite{dsp_proof}, the model first produces proof sketches, after which the automated prover Sledgehammer \citep{sledgehammer} is applied to complete the remaining gaps. As detailed in Table \ref{tab:base_math_pot_proof}, DeepSeekMath-Base 7B performs strongly on automated proof formalization.

\paragraph{Natural Language Understanding, Reasoning, and Code}

We analyze the model’s capabilities in natural language understanding with MMLU \citep{mmlu}, reasoning on BBH \citep{bbh}, and coding proficiency on HumanEval \citep{codex} and MBPP \citep{mbpp}. Table \ref{tab:understanding_reasoning_code} demonstrates that DeepSeekMath-Base 7B markedly improves results on MMLU and BBH relative to DeepSeek-Coder-Base-v1.5 \citep{deepseek-coder}, highlighting the advantages conferred by mathematical training for both understanding and reasoning in language tasks. Furthermore, by incorporating code-specific tokens in continued training, DeepSeekMath-Base 7B sustains the coding benchmark scores observed in DeepSeek-Coder-Base-v1.5. Taken together, DeepSeekMath-Base 7B consistently exceeds the performance of the generalist model Mistral 7B \citep{mistral} across all three evaluation categories for reasoning and coding.

\section{Supervised Fine-Tuning}

\subsection{SFT Data Curation}

We assemble a diverse instruction-tuning dataset in mathematics, encompassing both English and Chinese questions sourced from various mathematical domains and presenting a range of difficulty levels. Each question is paired with an explanatory answer in one of several formats: chain-of-thought (CoT) \citep{cot}, program-of-thought (PoT) \citep{pot,pal}, or tool-integrated reasoning \citep{tora}. Our final training corpus comprises 776K examples.

\begin{itemize}[topsep=0pt]
    \item \textbf{English mathematical datasets}: We augment GSM8K and MATH with tool-integrated solution annotations, utilize selected portions of MathInstruct \citep{MathInstruct}, and incorporate the Lila-OOD training data \citep{lila}, where solution formats include CoT or PoT. The English dataset features content from an array of mathematical disciplines, including algebra, probability, number theory, calculus, and geometry.
    \item \textbf{Chinese mathematical datasets}: Our Chinese corpus draws from K-12 mathematics, encompassing 76 sub-categories such as linear equations, and is annotated using both CoT and tool-integrated reasoning solution formats.
\end{itemize}

\subsection{Training and Evaluating DeepSeekMath-Instruct 7B}

Here, we describe the instruction tuning of DeepSeekMath-Instruct 7B, initialized from DeepSeekMath-Base. For training, samples are concatenated randomly until a sequence of up to 4K tokens is reached. The model is trained for 500 steps using a batch size of 256 and a constant learning rate of 5e-5.

Model evaluation targets mathematical reasoning capabilities, tested both with and without access to external tools, on four quantitative reasoning benchmarks spanning both English and Chinese. Performance is compared against state-of-the-art contemporaneous models:

\begin{itemize}[topsep=0pt]
    \item \textbf{Closed-source models} surveyed include: (1) members of the GPT series, with GPT-4 \citep{gpt4} and GPT-4 Code Interpreter~\footnote{\url{https://openai.com/blog/chatgpt-plugins\#\#code-interpreter}} representing the strongest versions, (2) Gemini Ultra and Pro \citep{gemini}, (3) Inflection-2 \citep{inflection-2}, (4) Grok-1~\footnote{\url{https://x.ai/model-card}}, as well as recent releases from Chinese technology firms, including (5) Baichuan-3~\footnote{\url{https://www.baichuan-ai.com}}, (6) the most recent GLM-4~\footnote{\url{https://open.bigmodel.cn/dev/api\#glm-4}}

    from the GLM family \citep{glm}.
\end{itemize}

The models discussed here are designed for general usage, with most having been subject to various alignment strategies.
    \item \textbf{Open-source models} encompass both broad-purpose and mathematically specialized architectures: (1) DeepSeek-LLM-Chat 67B \citep{deepseek-llm}, (2) Qwen 72B \citep{qwen}, (3) SeaLLM-v2 7B \citep{seallm}, and (4) ChatGLM3 6B \citep{chatglm3} serve as general models. Models with targeted mathematical abilities include (5) InternLM2-Math 20B~\footnote{\url{https://github.com/InternLM/InternLM-Math}}, which is based on InternLM2 with subsequent math-centric training and instruction tuning, (6) Math-Shepherd-Mistral 7B that incorporates PPO training \citep{schulman2017proximal} applied to Mistral 7B \citep{mistral} using a process-supervised reward model, (7) the WizardMath suite \citep{wizardmath}, which refines mathematical reasoning for Mistral 7B and Llama-2 70B \citep{llama2} through evolve-instruct (a variant of instruction tuning with AI-generated instructions) as well as PPO training on datasets primarily consisting of GSM8K and MATH problems, (8) MetaMath 70B \citep{metamath}, a Llama-2 70B variant fine-tuned with extended GSM8K and MATH datasets, (9) ToRA 34B \cite{tora}, derived from CodeLlama 34B and tuned for mathematical reasoning with tool integration, and (10) MAmmoTH 70B \citep{MathInstruct}, a Llama-2 70B model instruction-tuned using MathInstruct.
\end{itemize}

Table \ref{tab:sft_rl_math} displays results where, in a setup disallowing tool use, DeepSeekMath-Instruct 7B excels at performing stepwise reasoning. Remarkably, on the high-difficulty MATH benchmark, our model outperforms every open-source contender and most commercial models (such as Inflection-2 and Gemini Pro) by a margin of at least 9 percentage points, even outperforming models that are considerably larger (for instance, Qwen 72B) or have been subjected to advanced mathematics-oriented reinforcement learning (such as WizardMath-v1.1 7B). DeepSeekMath-Instruct is competitive with proprietary Chinese models like GLM-4 and Baichuan-3 on MATH, yet it does not match the leading results of GPT-4 or Gemini Ultra.

When the evaluation allows models to combine natural language reasoning with code or tool-based problem solving, DeepSeekMath-Instruct 7B achieves nearly 60\% accuracy on the MATH benchmark, outstripping all previously available open-source models. On other benchmarks, our model’s performance is on par with DeepSeek-LLM-Chat 67B, a much larger model regarded as the previous state-of-the-art.

\section{Reinforcement Learning}

\subsection{Group Relative Policy Optimization} 
Reinforcement learning (RL) has shown considerable effectiveness in further boosting the mathematical reasoning capacities of large language models following the Supervised Fine-Tuning (SFT) phase \citep{wang2023math, wizardmath}. Here, we present an efficient and impactful RL approach termed Group Relative Policy Optimization (GRPO).

\subsubsection{From PPO to GRPO}
Proximal Policy Optimization (PPO) \citep{schulman2017proximal} is a prominent actor-critic reinforcement learning method, widely adopted for RL fine-tuning in large language models \citep{ouyang2022training}. PPO seeks to improve model performance by maximizing the following surrogate objective:

\begin{equation}
\footnotesize
    \mathcal{J}_{PPO}(\theta) = \mathbb{E}{[q \sim P(Q), o \sim \pi_{\theta_{old}}(O|q)]} \frac{1}{|o|} \sum_{t=1}^{|o|} \min \left[ \frac{\pi_\theta(o_{t} | q, o_{<t})}{\pi_{\theta_{old}}(o_{t} | q, o_{<t})} A_{t}, \text{clip} \left( \frac{\pi_\theta(o_{t} | q, o_{<t})}{\pi_{\theta_{old}}(o_{t} | q, o_{<t})}, 1 - \

Here, $\pi_{\theta}$ represents the current policy while $\pi_{\theta_{old}}$ denotes the previous policy model; both are used to generate samples of questions $q$ and corresponding outputs $o$ from the dataset and $\pi_{\theta_{old}}$. The parameter $\varepsilon$ serves as a clipping hyper-parameter introduced by PPO to promote training stability. The advantage term, $A_t$, is estimated by applying Generalized Advantage Estimation (GAE) \citep{gae}, utilizing the rewards $\{r_{\ge t}\}$ in combination with a trainable value function $V_{\psi}$. As a result, PPO necessitates learning a separate value function alongside the policy network. To prevent the policy from excessively exploiting the reward model, it is standard practice to incorporate a token-level KL penalty against a reference policy in the reward signal for every token \citep{ouyang2022training}, i.e.,

\begin{equation}
    r_{t} = r_\varphi(q, o_{\le t}) - \beta \log\frac{\pi_{\theta}(o_{t}|q, o_{<t})}{\pi_{ref}(o_{t}|q, o_{<t})},
\label{eq:PPO-reward}
\end{equation}

where $r_\varphi$ corresponds to the learned reward model, $\pi_{ref}$ is the designated reference policy (commonly the pre-trained SFT model), and $\beta$ modulates the influence of the KL regularization.

Given that the value function model in PPO is usually sized comparably to the policy model itself, this approach entails significant demands on memory and computation. Moreover, the value function is used exclusively as a baseline in advantage calculations for the purpose of reducing variance. In large language model scenarios, typically only the final token receives a reward assignment, potentially complicating the effort to train a value function that yields accurate token-wise estimates. To overcome this, as visualized in Figure \ref{fig:grpo}, we introduce Group Relative Policy Optimization (GRPO). GRPO eliminates the necessity for learning an explicit value function as PPO does, and instead substitutes the group mean reward—computed across multiple sampled completions for the same input question—as the baseline. Specifically, for each prompt $q$, GRPO draws a set of completions $\{o_1, o_2, \cdots, o_G\}$ from $\pi_{\theta_{old}}$ and refines the policy by maximizing this objective:

\begin{equation}
\footnotesize
\begin{split}
    \mathcal{J}_{GRPO}(\theta) &= \mathbb{E}{[q \sim P(Q), \{o_i\}_{i=1}^G \sim \pi_{\theta_{old}}(O|q)]}  \\
    & \frac{1}{G}\sum_{i=1}^G\frac{1}{|o_i|} \sum_{t=1}^{|o_i|} \left\{ \min \left[ \frac{\pi_\theta(o_{i,t} | q, o_{i,<t})}{\pi_{\theta_{old}}(o_{i,t} | q, o_{i,<t})} \hat{A}_{i,t}, \text{clip} \left( \frac{\pi_\theta(o_{i,t} | q, o_{i,<t})}{\pi_{\theta_{old}}(o_{i,t} | q, o_{i,<t})}, 1 - \varepsilon, 1 + \varepsilon \right)  \hat{A}_{i,t} \right] - \beta \mathbb{D}_{KL}\left[\pi_{\theta} || \pi_{ref}\right]\right\} ,
\end{split}
\label{eq:GRPO-obj}
\end{equation}

Here, $\varepsilon$ and $\beta$ remain as adjustable hyper-parameters, and the advantage $\hat{A}_{i,t}$ is derived based solely on the relative rewards among outputs sampled for the same question, with further exposition to follow. GRPO’s mechanism of calculating advantages based on groupwise output comparisons is well-suited to the nature of reward modeling, as such models are often trained on datasets containing pairwise comparisons between candidate answers for a fixed prompt. It should also be emphasized that, as opposed to augmenting the reward with a KL penalty, GRPO directly incorporates a KL divergence regularizer between the candidate and reference policies into the loss. This simplification avoids complicating the evaluation of $\hat{A}_{i,t}$. Moreover, in contrast to the KL regularization approach of (\ref{eq:PPO-reward}), we use an unbiased estimator for the

\begin{algorithm}[t]
  \small
  \caption{Iterative Group Relative Policy Optimization}
  \textbf{Input} initial policy model $\pi_{\theta_{\text{init}}}$; reward models $r_\varphi$; task prompts $\mathcal{D}$; 
  hyperparameters $\varepsilon$, $\beta$, $\mu$
  \begin{algorithmic}[1]
    \State Initialize the policy model: $\pi_\theta \leftarrow \pi_{\theta_{\text{init}}}$
    \For{iteration = 1, \dots, I}
       \State Set the reference model $\pi_{ref} \leftarrow \pi_{\theta}$
      \For{step = 1, \dots, M}
      \State Draw a minibatch $\mathcal{D}_b$ from the dataset $\mathcal{D}$
      \State Assign the old policy model: $\pi_{\theta_{old}} \leftarrow \pi_{\theta}$ 
      \State For each question $q$ in $\mathcal{D}_b$, generate $G$ responses $\{o_i\}_{i=1}^G \sim \pi_{\theta_{old}} (\cdot \mid q)$
      \State For every sampled output $o_i$, evaluate rewards $\{r_i\}_{i=1}^{G}$ using $r_\varphi$
      \State Calculate group-relative advantages $\hat{A}_{i,t}$ for each token $t$ in $o_i$
      \For{GRPO iteration = 1, \dots, $\mu$}
        \State Update the policy $\pi_\theta$ by optimizing the GRPO objective as in Equation \ref{eq:GC-GRPO}
      \EndFor
    \EndFor
    \State Continuously train $r_\varphi$ via replay for ongoing reward model improvement.
    \EndFor
  \end{algorithmic}
  \textbf{Output} $\pi_\theta$
  \label{alg:iter-grpo}
\end{algorithm}

\subsubsection{Outcome Supervision RL with GRPO} 
Given a question $q$, the procedure samples a set of outputs $\{o_1, o_2, \cdots, o_G\}$ from the previous policy $\pi_{\theta_{old}}$. Each output is assessed by a reward model to produce rewards $\mathbf{r} = \{r_1, r_2, \cdots, r_G\}$. These scores are normalized within the group: subtracting the mean and dividing by the standard deviation. Outcome supervision employs the normalized reward at the end of each output $o_i$ as the advantage for every token in the output, meaning $\hat{A}_{i, t} = \widetilde{r}_i = \frac{r_i- {\rm mean}(\mathbf{r})}{{\rm std}(\mathbf{r})}$ for all $t$, and updates the policy by maximizing the objective in equation (\ref{eq:GRPO-obj}).

\subsubsection{Process Supervision RL with GRPO} 
While outcome supervision provides feedback only at the end of the sequence, such terminal reward signals may not be adequate for supervising complex multi-step reasoning tasks. Building on \cite{wang2023math}, we further apply process supervision, in which a reward is assigned following each intermediate reasoning step. For each question $q$ with $G$ outputs $\{o_1, o_2, \cdots, o_G\}$, a stepwise reward model scores each step, giving rewards: $\mathbf{R} = \{ \{r_1^{index(1)},\cdots,r_1^{index(K_1)}\}, \cdots,  \{r_G^{index(1)},\cdots,r_G^{index(K_G)}\} \}$. Here, $index(j)$ is the end position for the $j$-th step, and $K_i$ is the number of steps in $o_i$. The rewards are normalized across the batch: $\widetilde{r}_i^{index(j)} = \frac{r_i^{index(j)} - {\rm mean(\mathbf{R})}}{{\rm std(\mathbf{R})}}$.
Process supervision then determines the advantage at each token as the cumulative normalized reward over future steps: $\hat{A}_{i, t} = \sum_{index(j) \ge t} \widetilde{r}_i^{index(j)}$.
Policy optimization is carried out by maximizing the objective in equation (\ref{eq:GRPO-obj}).

\subsubsection{Iterative RL with GRPO} As reinforcement learning advances, earlier versions of the reward model may no longer provide adequate supervision for the evolving policy model. To address this, we investigate the iterative application of RL within the GRPO framework. In this approach, described in Algorithm \ref{alg:iter-grpo}, new reward model training datasets are constructed using samples generated by the latest policy model. The reward model is updated iteratively, leveraging a replay buffer that retains 10\% of prior training data to ensure continuity. Subsequently, the current policy model becomes the reference model, and policy optimization proceeds under the guidance of the updated reward model.

\subsection{Training and Evaluating DeepSeekMath-RL}

Our reinforcement learning experiments employ DeepSeekMath-Instruct 7B as the base model. The RL training dataset is composed of approximately 144K questions, formatted as chain-of-thought and sourced from GSM8K and MATH segments within the SFT dataset. We deliberately omit SFT data from other benchmarks to specifically analyze how RL impacts benchmarks for which no RL-phase data are provided. Following the procedure in \citep{wang2023math}, the reward model training data is prepared accordingly. The initial reward model is trained on DeepSeekMath-Base 7B with a learning rate of 2e-5. During GRPO, the policy model uses a learning rate of 1e-6, and the KL regularization coefficient is set to 0.04. For each problem, 64 outputs are sampled. The maximum output length is constrained to 1024 tokens, and the batch size for training is also set to 1024. The policy model undergoes only one update after every exploration cycle. Evaluation of DeepSeekMath-RL 7B adheres to the same benchmarks as used for DeepSeekMath-Instruct 7B. For DeepSeekMath-RL 7B, tasks from GSM8K and MATH with chain-of-thought annotations are classified as in-domain, whereas the remaining benchmarks are considered out-of-domain.

Table \ref{tab:sft_rl_math} presents the comparative results of open- and closed-source models employing both chain-of-thought and tool-augmented reasoning approaches on English and Chinese benchmarks. Our analysis indicates: 1) Utilizing chain-of-thought reasoning, DeepSeekMath-RL 7B achieves accuracies of 88.2\% on GSM8K and 51.7\% on MATH. These results exceed the performance of all open-source models within the 7B to 70B parameter range, and also outperform most closed-source counterparts. 2) Notably, DeepSeekMath-RL 7B is trained exclusively on instruction-tuning datasets in chain-of-thought format for GSM8K and MATH, initializing from DeepSeekMath-Instruct 7B. Even though its training dataset is relatively narrow in scope, DeepSeekMath-RL 7B achieves consistently better results across all evaluated metrics than DeepSeekMath-Instruct 7B, underlining the advantages conferred by reinforcement learning.

\section{Discussion}
In this section, we report observations and insights derived from our pre-training and RL experiments.

\subsection{Lessons Learnt in Pre-Training}

We begin by sharing our pre-training experience. Unless specified otherwise, all experiments adhere to the training configurations described in Section \ref{sec:quality-policy}. For clarity, throughout this section, references to the DeepSeekMath Corpus pertain to the 89B-token dataset generated during the second iteration of data collection.

\subsubsection{Code Training Benefits Mathematical Reasoning}

A frequently cited, though empirically underexplored, proposition is that code pre-training enhances reasoning skills. Our findings provide support for this claim in the context of mathematics: code training bolsters a model's proficiency in mathematical reasoning tasks, regardless of tool usage.

To evaluate the effect of code training on mathematical reasoning ability, we compared outcomes under two-stage training and one-stage training regimes:

\textbf{Two-Stage Training}

\begin{itemize}[topsep=0pt]
    \item \textbf{Code Training for 400B Tokens $\rightarrow$ Math Training for 150B Tokens}: DeepSeek-LLM 1.3B is first exposed to 400B tokens from code sources, after which it undergoes training on 150B tokens specifically curated for mathematics;
    \item \textbf{General Training for 400B Tokens $\rightarrow$ Math Training for 150B Tokens}: To provide a comparison, we run a control where, instead of code data, the model is pre-trained on 400B general-domain tokens (sampled from DeepSeek-AI's large-scale general corpus), before proceeding to the same 150B-token mathematics-focused stage. This helps examine whether code-specific data offers distinctive benefits over general textual data for the development of mathematical reasoning abilities.
\end{itemize}

\textbf{One-Stage Training}

\begin{itemize}[topsep=0pt]
    \item \textbf{Math Training for 150B Tokens}: The model DeepSeek-LLM 1.3B is directly trained on 150B math-specific tokens, without prior exposure to other domains;
    \item \textbf{Training on a mixture of 400B Code Tokens and 150B Math Tokens}: Because applying math training after an initial code pre-training leads to a drop in coding task performance, we further explore a single-phase training protocol in which code and math tokens (400B and 150B, respectively) are interleaved. This setup is used to determine whether such an approach can bolster mathematical reasoning while preserving code capabilities, thereby reducing catastrophic forgetting.
\end{itemize}

\paragraph{Results}
The comparative outcomes from these configurations are shown in Table \ref{tab:code-math} and Table \ref{tab:code-math-general-eval}.

Engaging in code token training consistently aids the model in program-assisted mathematical reasoning across both sequential (two-stage) and simultaneous (one-stage) paradigms. As evident from Table \ref{tab:code-math}, employing code data during the pre-training stage substantially enhances performance on GSM8K and MATH tasks via Python, even prior to targeted math token fine-tuning, with the latter further amplifying these gains. For the one-stage setting, co-training with both code and math data counters the loss in code proficiency that is typical after sequential fine-tuning and also yields a complementary effect in both programming (Table \ref{tab:code-math-general-eval}) and program-assisted mathematics (Table \ref{tab:code-math}).

Beyond program-based reasoning, code data also facilitates direct mathematical reasoning. Initial exposure to code in the two-stage paradigm already imparts notable benefits and enhances the efficacy of subsequent math-specialized training, culminating in optimal performance. Conversely, simultaneous code and math data training (one-stage) appears to diminish gains in mathematical reasoning when tools are not used. We hypothesize that the 1.3B parameter capacity of DeepSeek-LLM may not be sufficient to internalize both coding and advanced mathematics knowledge concurrently in a mixed training scenario.

\subsubsection{Limited Effectiveness of ArXiv Papers for Enhancing Mathematical Reasoning}

The inclusion of arXiv papers as a substantial component of mathematical pre-training corpora has become standard practice in recent large language model developments \citep{minerva,gpt-f,llemma,mathpile}. Nonetheless, in-depth examinations assessing their specific contribution to advancing mathematical reasoning capabilities are scarce. Contrary to common expectations, our empirical results suggest that arXiv papers do not meaningfully enhance mathematical reasoning. Our investigation employs models of varying scale, namely DeepSeek-LLM 1.3B and DeepSeek-Coder-Base-v1.5 7B \citep{deepseek-coder}, pre-trained on distinct arXiv-based corpora prepared using alternative processing strategies:

\begin{itemize}[topsep=0pt]
    \item \textbf{MathPile} \citep{mathpile}: consisting of 8.9B tokens curated using targeted cleaning and filtering heuristics, with scientific arXiv documents making up more than 85\% of the data;
    \item \textbf{ArXiv-RedPajama} \citep{redpajama}: comprising all arXiv LaTeX sources, stripped of preambles, comments, user macros, and bibliography entries, amounting to a total of 28.0B tokens.
\end{itemize}

We conduct isolated pre-training runs on both DeepSeek-LLM 1.3B (trained for 150B tokens) and DeepSeek-Coder-Base-v1.5 7B (trained for 40B tokens) using each of these arXiv-centric datasets. Our observations indicate that exposure to arXiv papers fails to facilitate measurable improvement in mathematical reasoning performance. When training exclusively on arXiv content, neither model demonstrates appreciable gains—and in some cases even shows regressions—across a spectrum of math benchmarks with varying levels of complexity considered in this work. The evaluated benchmarks include quantitative reasoning tasks (GSM8K, MATH, see Table \ref{tab:arxiv-cot}), multiple-choice STEM questions (MMLU-STEM, Table \ref{tab:arxiv-cot}), and formal mathematics assessments (miniF2F, Table \ref{tab:arxiv_atp}).

It is important to note that the scope of this conclusion is limited, and some important considerations remain unexamined:

\begin{itemize}[topsep=0pt]
    \item The potential utility of arXiv-based tokens for niche mathematical tasks excluded from our analysis, for instance, translating formal mathematical statements or proofs into informal language;
    \item The interactive effects resulting from combining arXiv data with other data types;
    \item Possible emergent benefits at larger model scales not considered in the present study.
\end{itemize}

We acknowledge these open questions and recommend that future work address them for a more comprehensive understanding.

\subsection{Reflections on Reinforcement Learning}
\subsubsection{Towards a Unified Paradigm} Here, we outline a unified conceptual framework for analyzing a range of training methodologies, encompassing SFT, RFT, DPO, PPO, and GRPO, and proceed with experiments to identify key determinants within this framework. In most cases, the gradient of the objective with respect to parameter $\theta$ in a given training method can be generally represented as follows:

\begin{equation}
    \nabla_{\theta}\mathcal{J}_{\textcolor{red}{\mathcal{A}}}(\theta) = \mathbb{E}[\underbrace{(q,o) \sim \textcolor{red}{\mathcal{D}}}_{Data \ Source}]\left( \frac{1}{|o|} \sum_{t=1}^{|o|}  \underbrace{GC_{{\mathcal{A}}}(q, o, t, \textcolor{red}{\pi_{{rf}}})}_{Gradient \ Coefficient}  \nabla_{\theta}\log \pi_{\theta}(o_t | q, o_{<t})\right).
\label{eq:objective}
\end{equation}

Three primary elements constitute the above framework: 1) the \textit{Data Source} $\mathcal{D}$, specifying the dataset utilized for training; 2) the \textit{Reward Function} $\pi_{{rf}}$, which supplies the evaluative feedback guiding optimization; and 3) the \textit{Algorithm} $\mathcal{A}$, responsible for integrating the data and reward to produce the gradient coefficient $GC$, governing the strength and direction of updates. The following commonly used techniques can be interpreted within this shared formulation:

\begin{itemize}
    \item  \textbf{Supervised Fine-tuning (SFT)}: In SFT, the pretrained model undergoes further training on datasets curated by humans specifically for this purpose.
    \item \textbf{Rejection Sampling Fine-tuning (RFT)}: RFT extends SFT by conducting additional fine-tuning using only those model outputs (for SFT-provided questions) which are retained after being evaluated for answer correctness.
    \item \textbf{Direct Preference Optimization (DPO)}: DPO refines the SFT model by optimizing it with a pair-wise preference loss over augmented samples obtained from the SFT model.
    \item \textbf{Online Rejection Sampling Fine-tuning (Online RFT)}: In contrast to standard RFT, Online RFT initializes with the SFT-trained model and subsequently improves it via fine-tuning on outputs generated by the actively updated policy model in real time.
    \item \textbf{PPO/GRPO}: With PPO/GRPO, the model is first initialized using the SFT checkpoint, after which policy improvement is performed using samples dynamically generated by the ongoing policy.
\end{itemize}

The elements of these approaches are outlined in Table \ref{tab:data_policy}. For a comprehensive explanation of the derivation, consult Appendix \ref{app:analysis-rl}.

\paragraph{Observation about Data Source} We categorize data sources as either online or offline sampling. Online sampling refers to obtaining training data from the real-time exploration trajectory of the current policy model during training, whereas offline sampling involves using data collected from the outputs of the initial SFT model. Both RFT and DPO utilize offline data, while Online RFT and GRPO operate with online-sampled data.

As illustrated in Figure \ref{fig:rl-analysis}, our results show that Online RFT considerably outperforms RFT on both evaluation benchmarks. Initially, Online RFT demonstrates performance similar to RFT; however, it establishes a clear lead as training progresses, showcasing the advantages of utilizing online training strategies. This phenomenon is expected: early in training, there is minimal distinction between the actor and the SFT model, resulting in sampled data that only exhibits slight variation. Over time, as the actor diverges from the SFT model, the sampled data reflects more pronounced differences, making the benefits of real-time data collection increasingly apparent.

\paragraph{Observation about Gradient Coefficient} The methodology incorporates the reward signal into the gradient coefficient used for model parameter updates. In our experimental setup, the reward function is sourced either from a `Rule'—evaluating response correctness—or from a `Model', in which a trained reward model assigns scores to responses. The reward model itself is supervised via rule-based assessment. Equations \ref{eq:GC-RFT} and \ref{eq:GC-GRPO} emphasize an important distinction: GRPO dynamically adjusts its gradient coefficient according to the output of the reward model, facilitating the proportional reinforcement or punishment of responses based on their reward magnitudes. By contrast, Online RFT lacks such adaptability, applying uniform positive reinforcement for all correct responses and omitting penalties for incorrect answers.

As illustrated in Figure \ref{fig:rl-analysis}, GRPO outperforms online RFT, underscoring the advantages of modifying positive and negative gradient coefficients. Additionally, GRPO+PS achieves better results than GRPO+OS, underscoring the value of employing step-specific, fine-grained gradient adjustments. We also investigate iterative RL, conducting two iterations in our experiments. According to Figure \ref{fig:iter-rl}, performance is markedly enhanced by iterative RL, particularly during the initial round.

\subsubsection{Why RL Works?}  
In this study, reinforcement learning is applied to a subset of the instruction tuning dataset, resulting in notable improvements over the baseline instruction-tuned model. To analyze the effectiveness of reinforcement learning, we examine the Pass@K and Maj@K metrics for both Instruct and RL-optimized models across two benchmarks. Figure \ref{fig:maj-pass} reveals that while RL improves the Maj@K metric, it does not impact Pass@K. This outcome suggests that RL enhances the general output distribution, i.e., \textbf{the primary improvement is linked to promoting correct answers within the TopK outputs, rather than bolstering the model’s core abilities.} In a similar vein, \citep{wang2023making} observed a \textbf{misalignment problem} in SFT models tackling reasoning tasks, and showed that reasoning skills could be improved by leveraging various preference alignment methodologies \citep{yuan2023rrhf,song2023preference,wang2023making}.

\subsubsection{How to Achieve More Effective RL?}
\label{sec:effec_rl}
Our findings confirm that RL delivers strong results for mathematical reasoning tasks. Furthermore, we present a unified framework to conceptualize a variety of notable training methods, interpreting them as either complete or reduced forms of RL approaches. This framework, as captured in Equation \ref{eq:objective}, incorporates three core elements: Data Source, Algorithm, and Reward Function. We also propose several avenues for future investigation regarding these elements.

\paragraph{Data Source} The data source constitutes the foundational input for all training approaches. Within the context of RL, we focus specifically on unlabeled questions paired with outputs sampled from the current policy. Our work utilizes questions stemming from the instruction tuning phase, and simple nucleus sampling to generate outputs. We attribute our RL pipeline’s selective improvement of Maj@K performance in part to this sampling strategy. For future research, we plan to extend our RL workflow to encompass out-of-distribution prompts and to integrate \textbf{advanced sampling (decoding) strategies}, such as tree-search-based techniques \citep{yao2023tree}. Additionally, \textbf{efficient inference techniques} \citep{xia-etal-2023-speculative,leviathan2023fast,kwon2023efficient,xia2024unlocking}—which dictate the exploration capability of policy models—are also anticipated to play a pivotal role.

\paragraph{Algorithms} The role of algorithms is to utilize both the collected data and the reward feedback to determine gradient coefficients, which subsequently adjust the model parameters. According to Equation \ref{eq:objective}, contemporary approaches almost universally place complete \textbf{TRUST} in the reward signal, leveraging it to modulate the conditional probability assigned to specific tokens, either enhancing or diminishing their likelihood. Nonetheless, this assumption of reward reliability frequently breaks down in practice, especially for highly intricate tasks. As a pertinent example, the PRM800K datasets \citep{lightman2023let}—despite being annotated by thoroughly trained individuals—still exhibit an error rate of around 20\%\footnote{\url{https://github.com/openai/prm800k/issues/12\#issuecomment-1728491852}}. Consequently, it is critical to investigate reinforcement learning algorithms that can tolerate or adapt to noisy reward feedback. Our position is that \textbf{WEAK-TO-STRONG} \citep{burns2023weak} alignment strategies could drive a significant transformation in the structure of learning algorithms.

\paragraph{Reward Function} The reward function fundamentally determines the feedback available for model training. In the context of RL, this function is most commonly instantiated as a neural reward model. We identify three pivotal research avenues for reward models: 1) \textbf{Improving the generalization capacity of the reward model.} It is imperative for the reward model to generalize effectively to out-of-distribution tasks and sophisticated decoding results, or else RL risks only maintaining the current output distribution of LLMs without real improvements in ability; 2) \textbf{Quantifying and expressing uncertainty in reward models.} Modeling uncertainty could facilitate interaction between less reliable (weak) reward signals and weak-to-strong learning methodologies; 3) \textbf{Constructing high-quality, process-level reward models with efficiency}, capable of delivering nuanced supervision to the underlying reasoning steps \citep{lightman2023let,wang2023math}.

\section{Conclusion, Limitation, and Future Work} 

This paper introduces DeepSeekMath, a model that sets a new benchmark among open-source systems on the competition-level MATH dataset, and reaches performance levels close to those of proprietary models. DeepSeekMath~builds on the DeepSeek-Coder-v1.5 7B foundation and is continually trained for 500B tokens, including a substantial subset of 120B math-specific tokens harvested from Common Crawl. Through a comprehensive ablation analysis, we find that web pages constitute a valuable reservoir for extracting high-quality mathematical content, whereas arXiv did not contribute as significantly as anticipated. We also propose Group Relative Policy Optimization (GRPO), an adaptation of Proximal Policy Optimization (PPO), which demonstrably enhances mathematical reasoning proficiency while maintaining low memory overhead. Empirical evidence suggests that GRPO offers notable gains, even as DeepSeekMath-Instruct 7B attains advanced benchmark results. Furthermore, we present a unified perspective to interpret a suite of algorithms and outline key avenues for further refining reinforcement learning methodologies.

Despite achieving strong results in quantitative reasoning tasks, DeepSeekMath~still shows comparatively weaker skills in areas like geometry and theorem-proving relative to leading closed-source systems. During preliminary evaluations, for example, the model struggled with triangle and ellipse-related questions, possibly reflecting bias introduced by data selection in both the pre-training and fine-tuning stages. Additionally, limitations in model scale mean that DeepSeekMath~does not match GPT-4's few-shot generalization; while GPT-4 benefits from few-shot examples, DeepSeekMath~exhibits nearly identical performance in both zero-shot and few-shot tests. Looking ahead, we plan to further optimize our data selection processes to create a richer and higher-quality training corpus. We will also continue to investigate directions described in Section \ref{sec:effec_rl} for advancing reinforcement learning techniques for LLMs.

\end{document}